\section{From $\opto_\varepsilon(K)$ oracles to the $\Hausdorf$ Problem}
\label{sec:haus}
% In this section, we show that we can get a result similar to that of~\Cref{upper-bound-1} while removing the well-separatedness condition. The well-separatedness assumption for a polytope $K$ is useful because for any vertex $v$ of $K$, there is a non-trivial probability of finding a direction $u$ along which the projection of $v$ is well-separated from that of other vertices of $K$. Hence, if we remove this assumption, we cannot hope to find a good approximation to each vertex of $K$ (for example, consider the case when $K$ closely approximates the unit circle on the two-dimensional plane). Instead we would like to approximate $K$ by another polytope $K'$, which may have many more vertices than $K$. We shall use the Hausdorff distance to measure the distance between two polytopes. 
We prove~\Cref{hausdorf-0} in this section. We begin by defining Hausdorff distance formally. 



\begin{defn}
\label{defn:Haus}
    The {\em Hausdorff-distance}, $\Haus(K,K')$, between two polytopes $K$ and $K'$ is the infimum over all values $\alpha$ such that the following condition is satisfied: for every point $x \in K$, there is a point $y \in K'$ such that $|x-y| \leq \alpha$, and vice versa.
\end{defn}

%
%We first restate~\Cref{Hausdorf-problem} in more precise terms:
We now give the formal version of~\Cref{hausdorf-0}:
\begin{theorem}\label{upper-bound-3}
Suppose $\varepsilon,\delta$ are reals in $[0,1]$ with 
\begin{align}
\label{eq:deltacond}
    \delta > c \varepsilon \sqrt{d}, \delta > c/\sqrt{d}, 
\end{align}
where $c$ is a large enough constant. 
  Let $K$ be a $k$-vertex polytope with $\Delta(K)$ denoting the diameter of $K$. Suppose we are also given an  $ \opto_\varepsilon(K)$ oracle $\cal O$.  
    Let 
  $P$ be the set of answers from the oracle on $m:= k^{10+c \delta^{-2}}$ independent random queries. Then, $\Haus(K, \CH(P)) \leq \delta \cdot \Delta(K)$. 
\end{theorem}


\begin{proof}
  We describe the algorithm for obtaining the desired set $S$ (referred as {\bf Random Probes Algorithm } in~\Cref{algorithmic-rsh})  in~\Cref{algo:S}. The set $P$ is constructed as follows: we pick a set  of $m$ random unit vectors. For each such unit vector $u$, we add the corresponding point $x(u)$ returned by the  oracle $\cal O$ to the set $P$. 

  \begin{algorithm}[H]
  \caption{Algorithm for finding the set $P$ such that $\CH(P)$ approximates $K$.}
  \label{algo:S}
    {\bf Input:} An $\opto_\varepsilon(K)$ oracle  $\cal O $. \;
        Initialize a set $P$ to $\emptyset$ \;
       {\bf Repeat} $m$ times:  \;
            \quad Let $u$ be a random unit vector in $\Re^d$. \;
            \quad Call $\cal O$ on $u$ to get a vector $x(u)$. \;
            \quad Add $x(u)$ to $P$. \;
        {\bf Output} $P$. 
\end{algorithm}

One side of the desired result is easy to show:
\begin{claim}
    \label{cl:Haus1}
    For each $x \in \CH(P)$, there is a $y \in K$ such that $|x-y| \leq \delta \Delta(K)$. 
\end{claim}
\begin{proof}
    For a point $x(u) \in P$, we know by the definition of $\opto_\varepsilon(K)$ oracle that there is a point $y \in K$ such that 
    $|x(u) - y| \leq \varepsilon \Delta(K) \leq \delta \Delta(K)$, where $\varepsilon \leq \delta$ follows from~\eqref{eq:deltacond}. The desired result now follows from the convexity of $K$. 
\end{proof}

It remains to show that for any vertex $v$ of $K$, there is a point $y \in \CH(P)$ such that $|v-y| \leq \delta \Delta(K)$. Fix such a vertex $v$ of $K$ for rest of the discussion. Let the random unit vectors considered  in~\Cref{algo:S} (in the order they get generated) be $u^1, \ldots, u^m$. Let $P^j$ denote the subset $\{x(u^1), \ldots, x(u^j)\}$ of $P$. Define an event $\cE_j$ as follows:
$$ \dist(v, \CH(P^j)) \leq \delta \cdot \Delta(K) \quad \mbox{or} \quad \dist(v, \CH(P^{j+1})) \leq \left(1 - \frac{\delta^2}{c'} \right) \dist(v, \CH(P^j)),$$
where $c'$ is a large enough constant. Our main technical result is to show that conditioned on {\em any} choice of $u^1, \ldots, u^j$ the event $\cE_j$ happens with reasonably high probability (where the probability is over the choice of $u^{j+1}$)$:$
\begin{lemma}
    \label{lem:tech1}
    For any index $j \in [m-1]$, 
    $$ \Pr_{u^{j+1}} \left[\cE_j|u^1, \ldots, u^j \right] \geq \frac{1}{100}. $$
\end{lemma}
\begin{proof}
    Fix the vectors $u^1, \ldots, u^j$. If $\dist(v, \CH(P^j)) \leq \delta \cdot \Delta(K)$, then we are done. So assume this is not the case. Let $b$ be the closest point in $\CH(P^j)$ to $v$. Thus, 
    \begin{align}
        \label{eq:b}
        |v-b| \geq \delta \cdot \Delta(K).
    \end{align}Define $w$ as $$ w:= \frac{v-b}{|v-b|}.$$
 We can now express the vector $u^{j+1}$ as $\lambda w + z$, where $\langle z, w \rangle = 0$. We first show the following useful properties of these vectors. 

 \begin{claim}
     \label{cl:tech2}
     With probability at least $\frac{1}{100}$, the following three events happen:
     \begin{align}
         |z| & \leq 4 \sqrt{d} \label{eq:event1} \\
         \max_{y \in K} |z \cdot (v-y)| & \leq 2 \sqrt{\ln k} \Delta(K) \label{eq:event2} \\
         \lambda & \geq \frac{100}{\delta} \sqrt{\ln k} \label{eq:event3}
     \end{align}
 \end{claim}
    \begin{proof}
        The proofs of these three inequalities are identical to those of~\Cref{fact:1},~\Cref{fact:2} and~\Cref{fact:3} (in order to prove~\eqref{eq:event2}, it suffices to show it for points $y$ which are vertices of $K$). 
    \end{proof}

The following fact is also easy to show:
\begin{fact}
    \label{fact:easy}
$$ |v-x(u^{j+1})| \leq 2 \Delta(K).$$
\end{fact}
\begin{proof}
    By the definition of $\opto_\varepsilon(K)$ oracle, there is a point $p \in K$ such that $|p-x(u^{j+1})| \leq \delta \cdot \Delta(K) \leq \Delta(K)$. The desired result now follows by triangle inequality. 
\end{proof}

 Let  $\delta_1$ denote $\frac{\delta^2}{100}$. Let $b_1$ denote the vector
 $$ \delta_1 x(u^{j+1}) + (1-\delta_1)b. $$
 Since $b_1 \in \CH(P^{j+1})$, the desired result will follow if we prove the following:
 \begin{align}
     \label{eq:tech4}
     |v-b_1|^2 \leq \left(1- \frac{\delta^2}{100} \right) |v-b|^2. 
 \end{align}

 Now, 
 \begin{align}
     \notag
     |v-b_1|^2 & = \delta_1^2 |v-x(u^{j+1})|^2 + (1-\delta_1)^2 |v-b|^2 + 2\delta_1(1-\delta_1) (v-x(u^{j+1}))\cdot (v-b) \\
    \notag & \stackrel{\mbox{\small{\Cref{fact:easy}}}}{\leq} 4 \delta_1^2 \Delta(K)^2 + (1-2\delta_1) |v-b|^2 + \delta_1^2 \Delta(K)^2 + 2\delta_1(1-\delta_1) (v-x(u^{j+1}))\cdot (v-b) \\
    \notag &  \stackrel{\eqref{eq:b}}{\leq} \left( 1 - \frac{3\delta_1}{2} \right) |v-b|^2 +  + 2\delta_1(1-\delta_1) |v-b| \cdot (v-x(u^{j+1}))\cdot w \\
    \notag & = \left( 1 - \frac{3\delta_1}{2} \right) |v-b|^2 +  + 2\delta_1(1-\delta_1) \left( \underbrace{\frac{ |v-b|}{\lambda} \cdot (v-x(u^{j+1}))\cdot u^{j+1}}_{:=A}  \right. \\ \label{eq:AB}
  & \quad \quad  \left. \underbrace{-\frac{ |v-b| }{\lambda}\cdot (v-x(u^{j+1}))\cdot z} _{:=B} \right)
 \end{align}
 We now bound each of the terms $A$ and $B$ above. Now, 
 $$ A \leq \frac{|v-b|}{\lambda} \varepsilon \Delta(K) \stackrel{\eqref{eq:deltacond}}{\leq} \frac{|v-b| \delta \Delta(K)}{c \lambda} \stackrel{\eqref{eq:b}}{\leq} \frac{|v-b|^2}{c \lambda} \stackrel{\eqref{eq:event3}}{\leq} \frac{|v-b|^2 \delta}{c},$$
 where the first inequality follows from the definition of $\opto_\varepsilon(K)$ oracle. We now bound the quantity $B$. Let $y$ be the point in $K$ closest to $x(u^{j+1})$. We know that $|y-x(u^{j+1})| \leq \varepsilon \Delta(K)$. Therefore, 
 \begin{align*}
     B & \leq \frac{|v-b|}{\lambda} \left( |(v-y) \cdot z| + |z| \varepsilon \Delta(K) \right) \\
     & \stackrel{\eqref{eq:event1}, \eqref{eq:event2}}{\leq} \frac{|v-b|}{\lambda} \left( 2 \sqrt{\ln k} \Delta(K) + 4 \varepsilon \sqrt{d} \Delta(K) \right) \\
     & \stackrel{\eqref{eq:deltacond}, \eqref{eq:event3}}{\leq} \frac{\delta |v-b| \Delta(K) }{10} \stackrel{\eqref{eq:b}}{\leq} \frac{|v-b|^2}{10}. 
 \end{align*} 
 Substituting the above bound on $A$ and $B$ in~\eqref{eq:AB} yields the desired result. 
\end{proof}
We are now almost done. As the following result shows, it suffices to argue that enough number of events $\cE_j$ happen:
\begin{claim}
\label{cl:close}
    If at least $\frac{c'}{\delta^2} \ln(2/\varepsilon)$ of the events $\cE_j, j \in [m-1]$ happen, then $\dist(v, \CH(P)) \leq \delta \Delta(K)$. 
\end{claim}
\begin{proof}
Assume, for the sake of contradiction, that $\dist(v,\CH(P)) > \varepsilon \Delta(K)$. 
Assume that events $\cE_{j_1}, \ldots, \cE_{j_h}$ happen, where $h :=\frac{c'}{\delta^2} \ln(2/\varepsilon). $ Now, for any index $i \in [h-1]$, the definition of $\cE_{j_{i+1}}$ implies that 
$$\dist(v, \CH(P^{j_{i+1}})) \leq \left( 1 - \frac{\delta^2}{c'} \right) \dist(v, \CH(P^{j_{i+1}-1}) \leq  \dist(v, \CH(P^{j_{i}})).$$
Therefore, 
$$\dist(v, \CH(P^{j_h})) \leq \left( 1 - \frac{\delta^2}{c'} \right)^h \dist(v, P^1) 
\stackrel{\mbox{\small{\Cref{fact:easy}}}}{\leq} \left( 1 - \frac{\delta^2}{c'} \right)^h 2 \Delta(K) \leq \delta \Delta(K). $$
\end{proof}

It remains to show that with high probability at least $h :=\frac{c'}{\delta^2} \ln(2/\varepsilon) $
of the events happen. In order to prove this, we divide the sequence $[m]$ into $[h]$ subsequences, each of length $m/h$. Call these subsequences $C_1, \ldots, C_h$. It follows from~\Cref{lem:tech1} that for any $i \in [h]$, 
$$ \Pr \left[ \wedge_{j \in C_i} \neg \cE_j \right] \leq 0.99^{m/h} \leq \frac{1}{h^2}.$$
A simple union bound now shows that with probability at least $1-1/h$, at least one event $\cE_j$ happens during each of the subsequences $C_1, \ldots, C_h$.~\Cref{cl:close} now proves the theorem. 
% First, we relate the parameters here to the parameters we had in our write-up. Under ``Assumptions on model and generation process'', (A1): Mean Separation, we had (with Diameter $\leq 1$ and $\Gamma $ replaced by $\delta$, the inequality we had connecting parameters (it was (18)) in the Dec 16 old write-up):
% \begin{equation}\label{791}
% (i) \delta \in \Omega(\sqrt k\sigma_0/\sqrt{w_0})\quad ;\quad (ii) \delta \in \Omega^*(1).\end{equation}
% Now, the error $\varepsilon$ is bounded above by $\sigma_0/\sqrt{w_0}$ in our old write-up since for a set $S$ of $w_0n$ elements, we have $|A_{\cdot,S}-P_{\cdot,S}|\leq \sigma_0/\sqrt{w_0}$. We will see here (using Thm 3.1 above in the sVD projection) that $\varepsilon\in \Omega(1/\sqrt k)$ suffices. To get this we need 
% $$\sigma_0\in O(\sqrt{w_0}/\sqrt k)$$ which is exactly inequality (i) in (\ref{791}) since $\delta \in\Omega(1)$.

% We proved already that if hat denotes projection to SVD subspace, (Lemma 4.5)
% $$|\widehat M_{\cdot,\ell}-M_{\cdot,\ell}|\leq \sigma_0/\sqrt {w_0}.$$
% Thus, $\delta $ is changed by at most a factor of 2 when we project.

% Note that by subset smoothing, we have an optimization oracle in ${\bf R}^d$ with error $\varepsilon\leq \sigma_0/\sqrt{w_0}$ which is at most $O(1/\sqrt k)$.
% Now for any $u$ in the SVD subspace, the same oracle automatically serves as an opt oracle over the projection with at most the same error.

\end{proof}
